\section{Joint tests and survival of field labels}
\label{sec:boundary-detection}
Two details complete the common-domain reconstruction, apart from locality
and the momentum projection measure. First, the joint tempered distributions
supply the test-function continuity used in the word construction. Second,
a label detectable by Euclidean correlations remains detectable after
continuation to real time. The latter uses the full-tube construction and
its identification with the ordered-time function; no pointwise boundary
values of the fields are required.

\begin{lemma}\label{lem:tensor-tests}
The finite linear span of functions
$\prod_{i=1}^n f_i(X_i)$, $f_i\in\mathscr S(\mathbb R^4)$, is dense in
$\mathscr S((\mathbb R^4)^n)$. The product map from the $n$ test spaces
into the joint test space is continuous multilinear.
\end{lemma}
\begin{proof}
Every joint derivative of a product separates into the corresponding
single-slot derivatives. Also
$1+|X|\leq\prod_i(1+|X_i|)$ up to a fixed constant depending on $n$.
Thus each weighted derivative supremum of the product is bounded by a
constant times a product of single-slot Schwartz seminorms. This proves
continuous multilinearity.

For density, compactly supported smooth joint tests are already dense by
the cutoff argument in Proposition~\ref{prop:boost-ward}. Fix such a test
$f$, and choose single-slot smooth compact cutoffs $\chi_i$ with product
$\chi(X)=\prod_i\chi_i(X_i)$ equal to one on $\operatorname{supp}f$.
Let $\eta_R=\eta(\cdot/R)$, where $\eta$ is a smooth compactly supported
frequency cutoff equal to one near zero. By the cutoff argument of
Proposition~\ref{prop:boost-ward}, applied in frequency space,
$\eta_R\widehat f\to\widehat f$ in Schwartz space as $R\to\infty$.
Since Fourier inversion is continuous on Schwartz space and multiplication
by the fixed function $\chi\in C_c^\infty$ is continuous there,
\[
 f_R=\chi\,\mathcal F^{-1}(\eta_R\widehat f)\longrightarrow\chi f=f
\]
in Schwartz space. Here the Fourier transform and its inverse are continuous
inverse maps on that space, with the convention of
Section~\ref{sec:fourier-support}; see \cite[Theorem 1.3, p.22]{Melrose}.
Multiplication by the fixed compact cutoff is continuous by the Leibniz rule.

For fixed $R$, put $g=\eta_R\widehat f$. Fourier inversion writes
\[
 f_R(X)=(2\pi)^{-4n}\int_{\operatorname{supp}g}
       g(p)\prod_i\big[\chi_i(X_i)e^{ip_i\cdot X_i}\big],dp.
\]
This integral is a limit in Schwartz space of finite sums of its displayed
product tests. To verify that assertion explicitly, enclose the compact
frequency support in a box and use ordinary rectangular Riemann sums.
All summands and the integral have the fixed compact position support
$\operatorname{supp}\chi$. On the product of that support with the
frequency box, every position derivative of the integrand is continuous,
hence uniformly continuous. The error in each weighted derivative
supremum is bounded by the frequency-box volume times its uniform modulus
of continuity at the partition diameter, times the maximum position weight
on the compact support. It tends to zero. Thus the same sequence of
partitions converges in every Schwartz seminorm. Each Riemann sum belongs
to the stated span, so does $f_R$ in its closure, and taking $R\to\infty$
proves density. For $n=0$ the joint space is one-dimensional and the empty
product one already spans it.
\end{proof}

\begin{corollary}\label{cor:word-test-inputs}
For the jointly tempered distributions constructed in
Theorem~\ref{thm:joint-boundary}, the word functional is separately linear
in every label and test. Its mixed Gram expressions
\[
 (f,g)\longmapsto K((v,f)b,(u,g)c)
\]
are jointly continuous. Thus these parts of
Input~\ref{input:word-boundary} follow from the construction.
\end{corollary}
\begin{proof}
Multilinearity follows by composing the label-multilinear distribution
with the multilinear product test. In the Gram expression the first
variable occurs as $\overline f$ in a fixed slot and the second as $g$
in another fixed slot; all other factors and labels are fixed. Conjugation
is continuous on Schwartz space. Lemma~\ref{lem:tensor-tests} and continuity
of the fixed-label distribution therefore give the asserted joint
continuity. The empty-word value is one because the zero-insertion
continuation is constantly one.
\end{proof}

The remaining hypothesis \eqref{eq:word-motion-input} of
Proposition~\ref{prop:word-motion} concerns the Poincar\'e group. Let
$g=(a,\Lambda)$ act on spacetime by $gX=\Lambda X+a$, with $\Lambda$
proper orthochronous, and on tests by $P_gf=f\circ g^{-1}$, as in
Lemma~\ref{lem:schwartz-motion}. Let $\rho$ be the label action of
Proposition~\ref{prop:lorentz-label-descent}, on which translations act
trivially. Define
\[
 g(v,f)=(\rho(\Lambda)v,P_gf)
\]
on letters, extend it letter by letter to words, and let it fix the empty
word. The word construction uses the label involution $J$ of
Section~\ref{ex-reflection}, the reflected conjugate. The following
corollary assumes that this is the involution with which
Proposition~\ref{prop:lorentz-label-descent} proves commutation, that is,
the involution of Input~\ref{input:calibrated-actions}.

\begin{corollary}[Poincar\'e invariance of the word kernel]
\label{cor:word-motion-inputs}
Assume the hypotheses of Theorem~\ref{thm:lorentz-covariance} for the
action $\rho$ just described, and assume the identification of involutions
just stated. Then $b\mapsto gb$ is an action of the Poincar\'e group on
words, $(gb)^*=g(b^*)$, and
\[
 K(gb,gc)=K(b,c),\qquad g\longmapsto K(b,gc)\ \text{is continuous},
\]
for all words $b,c$. Thus \eqref{eq:word-motion-input} holds. The
translation hypothesis of Proposition~\ref{prop:word-motion} is supplied,
for every $t>0$, by Proposition~\ref{prop:heat-cups}.
\end{corollary}
\begin{proof}
For $g,h$ in the Poincar\'e group,
$P_g(P_hf)(X)=f(h^{-1}g^{-1}X)=P_{gh}f(X)$, the Lorentz part of $gh$ is
the product of the Lorentz parts, and $\rho$ is a representation. Hence
the letterwise rule is a group action on words, and $P_e$ and $\rho(I)$
are identities.

Proposition~\ref{prop:lorentz-label-descent} gives $J\rho(\Lambda)=\rho(\Lambda)J$.
Since $P_g$ commutes with complex conjugation of tests,
\[
 (g(v,f))^*=(J\rho(\Lambda)v,\overline{P_gf})
          =(\rho(\Lambda)Jv,P_g\overline f)=g\big((v,f)^*\big).
\]
The star of a word reverses the order of its letters and stars each letter.
Reversal commutes with a letterwise action, so $(gb)^*=g(b^*)$ and
$(gb)^*(gc)=g(b^*c)$.

Now let $d=(v_1,f_1)\cdots(v_n,f_n)$ be any word and
$F=\prod_if_i(X_i)$ its product test. The product of the moved tests is
$\prod_i(P_gf_i)(X_i)=F(g^{-1}X_1,\ldots,g^{-1}X_n)$, the simultaneous
pullback of $F$ by $g$. Write $g=(a,I)(0,\Lambda)$. The translation
invariance and Lorentz covariance in Theorem~\ref{thm:lorentz-covariance}
give
\[
 w(gd)=W_n(\rho(\Lambda)v)\big(P_{(a,I)}P_{(0,\Lambda)}F\big)
      =W_n(\rho(\Lambda)v)\big(P_{(0,\Lambda)}F\big)=W_n(v)(F)=w(d).
\]
Applying this to $d=b^*c$ gives $K(gb,gc)=w(g(b^*c))=w(b^*c)=K(b,c)$.

For continuity fix words $b$ of length $n$ and
$c=(u_1,\varphi_1)\cdots(u_m,\varphi_m)$. All labels lie in one finite level $E_d$.
Choose a basis $e_1,\ldots,e_N$ of $E_d$ and write
$\rho(\Lambda)u_j=\sum_kr_{kj}(\Lambda)e_k$. The coefficients $r_{kj}$ are
continuous, because $\rho$ is continuous on each finite level. Let
$F_g$ be the joint test of $b^*(gc)$. It equals $F_e\circ\widetilde g^{-1}$,
where $\widetilde g$ is the invertible affine map of the joint position
space which acts by the identity on the first $n$ positions and by $g$ on
each of the last $m$ positions. Multilinearity of $W_{n+m}$ in the labels
gives
\[
 K(b,gc)=\sum_{k_1,\ldots,k_m}\Big(\prod_{j=1}^mr_{k_jj}(\Lambda)\Big)
   W_{n+m}(Jv_n,\ldots,Jv_1,e_{k_1},\ldots,e_{k_m})(F_g),
\]
where $v_1,\ldots,v_n$ are the labels of $b$. Each fixed-label
distribution is continuous on the joint Schwartz space. The map
$g\mapsto\widetilde g$ is continuous into the group of invertible affine
maps of the joint position space, so the general form of
Lemma~\ref{lem:schwartz-motion} shows that $g\mapsto F_g$ is continuous
into the joint Schwartz space. The finite
sum of products of continuous functions is continuous in $g$.

The last statement is Proposition~\ref{prop:heat-cups}: for $g=(a,I)$
the word action translates every test, and the heat vectors satisfy
$h_t(\tau_ab)=T_E(a)h_t(b)$ for every $t>0$.
\end{proof}

The identification of involutions is a same-data requirement, not a
convention. Let $\theta$ be the label action of the time reflection and
$\kappa$ the real structure, so that the reflected conjugate is
$J=\theta\kappa$. The time reflection conjugates the rotation in the
time--$a$ plane to its inverse, so $\theta L_a\theta=-L_a$. If the
Euclidean generators are real, $\kappa L_a\kappa=L_a$, antilinearity of
$J$ gives
\[
 J(-iL_a)J^{-1}=i\,\theta L_a\theta=-iL_a,
\]
so $J$ commutes with the boost generators $K_a=-iL_a$, while $\kappa$ alone
would anticommute with them. Both $\theta$ and $\kappa$ commute with spatial
rotations.

\begin{lemma}[A zero boundary has zero strict continuation]
\label{lem:zero-boundary}
Fix an ordered label tuple $v$. If $W_n(v)=0$ as a joint tempered
distribution, then $F_n(t+iy,x;v)=0$ for every strictly ordered height tuple
$y$ and every real position tuple $(t,x)$.
\end{lemma}
\begin{proof}
The full-tube function is defined by
$H_n(v;z)=(2\pi)^{-4n}\widehat W_n(v)(L_z)$.
It is therefore identically zero if $W_n(v)=0$. Proposition~\ref{prop:tube-identification}
identifies its restriction at $z=X+i d(y)$ with the ordered-time function
for every strictly ordered $y$. This proves the assertion.
\end{proof}

Let $T$ denote the selected Euclidean hierarchy. At each finite raw label
level, its null space consists of labels whose insertion vanishes in every
distinct-point correlation with arbitrary exterior labels. Assume the
selected and original null spaces agree at every level, as established by
the detector argument when its moment and rank hypotheses are supplied.
Assume also the finite-level quotient embeddings into the union are
injective and every union label has a finite-level representative.

\begin{proposition}[Boundary detection]\label{prop:boundary-detection}
If a union label $v$ satisfies
\[
 w(b(v,f)c)=0
\]
for every pair of exterior words $b,c$ and every single-slot Schwartz test
$f$, then $v=0$. Consequently every nonzero label defines a nonzero smeared
field operator in Proposition~\ref{prop:word-fields}.
\end{proposition}
\begin{proof}
Take any ordered tuple of labels containing $v$, at any slot. The
assumption says that its joint distribution vanishes on every product
test: group the factors before the distinguished slot into $b$, and those
after it into $c$. Lemma~\ref{lem:tensor-tests} and tempered continuity
therefore make that entire joint distribution zero. By
Lemma~\ref{lem:zero-boundary}, its strict-time continuation is zero.
At real times zero, the pure-imaginary identity gives zero for the
corresponding selected Euclidean correlation whenever the Euclidean times
are strictly increasing.

For arbitrary positions with pairwise distinct time coordinates, sort
them by time. The label tuple is permuted with the positions and still
contains $v$, so the preceding zero-distribution argument applies to the
sorted tuple. Euclidean permutation symmetry transfers its zero value
back to the original ordering. Such position tuples are dense in the
full configuration space of distinct points: perturb the finitely many
time coordinates by distinct arbitrarily small amounts, avoiding the
finitely many equality hyperplanes. Openness of the distinct-point
configuration space keeps all points distinct during a sufficiently small
perturbation. Continuity of the Euclidean correlations extends the zero
value to every configuration, including those with time ties.

Represent $v$ by a finite-level raw label. The preceding conclusion,
with arbitrary raw exterior labels projected into the union, is exactly
the all-exterior null criterion for that label in $T$. Equality of selected
and original null spaces makes its original quotient class zero. Its
union image is $v$, so $v=0$. This proves the detection premise of
Proposition~\ref{prop:word-locality}; that proposition then supplies the
nonzero smeared operator for every nonzero label.
\end{proof}

This argument closes the boundary-detection step conditional on the already
specified equality of Euclidean null spaces. The estimate and finite-rank
inputs establishing that equality remain part of the quantitative phase
construction.
